\section*{Hoofdstuk \ref{ch:b2:acyl-substitution} --- Carbonzuurderivaten}

\begin{solution}{exo:b2:acyl-substitution:1}
Ethaanamide $<$ ethylethanoaat $<$ ethaanzuuranhydride $<$ ethanoylchloride.
\end{solution}

\begin{solution}{exo:b2:acyl-substitution:2}
Ethaanzuur (+ \ce{HCl}); ethylethanoaat; ethaanamide (+ ammoniumchloride, waarbij het tweede equivalent
ammoniak het \ce{HCl} wegvangt); ethaanzuuranhydride (+ \ce{NaCl});
\textit{N},\textit{N}-dimethylethaanamide (+ triethylammoniumchloride).
\end{solution}

\begin{solution}{exo:b2:acyl-substitution:3}
\ce{C6H5COOCH3 + NaOH -> C6H5COONa + CH3OH}. $13.6/136.15 = \qty{0.0999}{mol}$, evenveel
\ce{NaOH}: \qty{4.00}{g}.
\end{solution}

\begin{solution}{exo:b2:acyl-substitution:4}
$\gamma$-Butyrolacton (oxolaan-2-on): een vijfring van vier koolstofatomen en één zuurstofatoom.
\end{solution}

\begin{solution}{exo:b2:acyl-substitution:5}
Ammoniak addeert aan het carbonylkoolstofatoom; een proton gaat van het stikstofatoom naar het zuurstofatoom (of naar de
vertrekkende groep); ethoxide (als ethanol, na protonoverdracht) vertrekt en de \ce{C=O} vormt zich opnieuw:
ethaanamide en ethanol.
\end{solution}

\begin{solution}{exo:b2:acyl-substitution:6}
De fenolische \ce{OH} wordt geacyleerd (zure katalyse): 2-acetoxybenzoëzuur; bijproduct ethaanzuur.
\end{solution}

\begin{solution}{exo:b2:acyl-substitution:7}
2-Methylpropaan-2-ol. De eerste additie geeft propanon, elektrofieler dan de ester, dat
meteen de tweede methylgroep opneemt.
\end{solution}

\begin{solution}{exo:b2:acyl-substitution:8}
Butaan-1,4-diol: de ester wordt gereduceerd tot twee primaire alcoholen, waarbij de ring opengaat.
\end{solution}

\begin{solution}{exo:b2:acyl-substitution:9}
Eén equivalent: $x^2/(1 - x)^2 = 4$, $x = 2/3$. Drie equivalenten: $x^2/[(1 - x)(3 - x)] = 4$, $3x^2 - 16x
+ 12 = 0$, $x = 0.90$. Het water verwijderen naarmate het ontstaat (Dean--Stark) drijft de reactie verder.
\end{solution}

\begin{solution}{exo:b2:acyl-substitution:10}
\ce{SOCl2} geeft benzoylchloride; met twee equivalenten diethylamine (of één plus triethylamine)
geeft dat het amide. Zuur en amine die rechtstreeks gemengd worden, ondergaan een zuur-basereactie: diethylammoniumbenzoaat,
dat pas bij sterk verhitten het amide geeft.
\end{solution}

\begin{solution}{exo:b2:acyl-substitution:11}
Drie KOH per trioleïne: $3 \times 56.11/885.4 = \qty{0.190}{g}$ per gram, een verzepingsgetal van
190.
\end{solution}

\begin{solution}{exo:b2:acyl-substitution:12}
\ce{NaBH4}: alleen het keton, tot een secundaire alcohol. \ce{LiAlH4}: het keton tot de secundaire
alcohol, de ester tot een primaire alcohol (waarbij het alcoholdeel vrijkomt), het amide tot een amine.
\end{solution}

\begin{solution}{pb:b2:acyl-substitution:1}
\textbf{1.} Oliezuur \ce{C18H34O2}; glycerol \ce{C3H8O3}.
\textbf{2.} \ce{C57H104O6}: $57 \times 12.011 + 104 \times 1.008 + 6 \times 15.999 = \qty{885.45}{g/mol}$.
\textbf{3.} Drie.
\textbf{4.} De cis-dubbele binding geeft de ketens een knik, zodat ze slecht stapelen: zwakke intermoleculaire krachten, een
laag smeltpunt.
\textbf{5.} Glycerol en drie natriumoleaat, een zeep.
\textbf{6.} Een mengsel van methylesters van vetzuren.
\textbf{7.} \ce{C57H104O6 + 3CH3OH -> C3H8O3 + 3C19H36O2}.
\textbf{8.} Het methoxide-ion, gevormd uit methanol en de base, is een veel sterker nucleofiel dan
methanol; het valt de carbonylgroepen van de esters aan.
\textbf{9.} Het carbonylkoolstofatoom van één estergroep, dat \ce{O-}, de \ce{OCH3} van het methoxide en het
glycerylzuurstofatoom draagt: tetraëdrisch.
\textbf{10.} \omterm{def:b2:acyl-substitution:transesterification}{Omestering} is een evenwicht met een constante rond 1; een overmaat methanol verschuift het
naar de methylesters.
\textbf{11.} Glycerol scheidt zich af als onderste laag: een product verwijderen drijft het evenwicht voort.
\textbf{12.} Nee: methoxide wordt bij elke uitwisseling teruggevormd.
\textbf{13.} Het neutraliseert de base en geeft natriumoleaat (zeep): de katalysator wordt verbruikt.
\textbf{14.} Het verzeept ze: hydroxide (uit water en methoxide) knipt de esters onomkeerbaar door tot zeep en
methanol.
\textbf{15.} Ze verbruikt base, emulgeert het mengsel en verhindert dat de glycerollaag zich afscheidt.
\textbf{16.} Door eerst een zuurgekatalyseerde \omterm{def:b2:acyl-substitution:esterification}{verestering} met methanol, die de vrije zuren omzet in
methylesters.
\textbf{17.} Water hydrolyseert esters en zet de katalysator om in hydroxide, dat verzeept.
\textbf{18.} $10^6/885.45 = \qty{1129}{mol}$.
\textbf{19.} $3 \times 1129 \times 32.04 = \qty{108.6}{kg}$.
\textbf{20.} $3 \times 1129 \times 296.50 = \qty{1004.6}{kg}$.
\textbf{21.} $1000 + 108.6 = \qty{1108.6}{kg}$ erin; $104.0 + 1004.6 = \qty{1108.6}{kg}$ eruit.
\textbf{22.} \qty{20}{kg} oliezuur, $20\,000/282.47 = \qty{70.8}{mol}$, evenveel \ce{NaOH}:
\qty{2.83}{kg}.
\textbf{23.} $1129 \times 92.09 = \boldsymbol{\approx \qty{104}{kg}}$ glycerol per ton.
\end{solution}
